\documentclass[a4paper]{article}

\usepackage{graphicx}
\usepackage{amsmath,amssymb}
\usepackage{minted}
\usepackage{multirow}
\usepackage{float}
\usepackage{todonotes}
\usepackage{forest}
\usepackage{xstring}
\usepackage{xcolor}
\usepackage[most]{tcolorbox}
\usepackage{xparse}
\usepackage{natbib}

\usetikzlibrary{graphs,graphdrawing}
\usegdlibrary{layered}

% for external graphviz
\input{/home/donkarlo/Dropbox/repo/nd_ai_project/src/nd_ai/knoweldge_representation/language/formal/computer/programming/tex/latex/code_clips/external_graphviz.tex}

% for tags
\input{/home/donkarlo/Dropbox/repo/nd_ai_project/src/nd_ai/knoweldge_representation/language/formal/computer/programming/tex/latex/parts/control_sequence/kind/semtag/semtag.tex}

% for links
\input{/home/donkarlo/Dropbox/repo/nd_ai_project/src/nd_ai/knoweldge_representation/language/formal/computer/programming/tex/latex/code_clips/seehere.tex}

\title{Default mode network}
\date{}

\begin{document}
    \maketitle
    Being active when a person is not focused on the outside world and the brain is at wakeful rest, such as during daydreaming and mind-wandering \seeherefooter{https://en.wikipedia.org/wiki/Default_mode_network}
    
    \cite{raichle-2015-the-brains-default-mode-network} reviews the discovery and development of the concept of the brain’s default mode network. The paper explains that the DMN was first noticed through task-induced decreases in activity during attention-demanding tasks.
    It describes the DMN as a set of bilateral cortical regions involving medial prefrontal, posterior cingulate/precuneus, parietal, and temporal areas. The review connects the DMN to resting-state functional connectivity, intrinsic brain activity, memory, self-related processing, and emotional regulation. Raichle argues that the DMN is not merely a mind-wandering network, but a central part of the brain’s intrinsic organization in health and disease.
    
    
    \section{Relation with reviewing memories}
        \cite{andrews-hanna-2010-evidence-for-the-default-networks-role-in-spontaneous-cognition}
        examine whether the default mode network supports spontaneous cognition. They show that activity in default-network regions increases when people report internally directed thoughts. These thoughts often include autobiographical memories, future simulations, and self-related mental content. The paper argues that the default network is not simply “resting,” but actively supports mental processes detached from the external task. Its main contribution is linking DMN activity to spontaneous thought, especially personal past and future thinking.
    
    
    \section{Relation with internal narrative or internal language}
        \cite{menon-2023-20-years-of-the-default-mode-network-a-review-and-synthesis} reviews two decades of research on the Default Mode Network and explains how its role has expanded from a “resting-state” network to a central system for self-reference, social cognition, episodic memory, semantic memory, language, and mind wandering. The paper argues that the Default Mode Network is not a single-purpose system, but a set of interacting nodes and subnetworks whose functions depend on dynamic interactions with salience and frontoparietal control networks. Its main synthesis is that the Default Mode Network integrates memory, language, and semantic representations to construct a coherent internal narrative, which supports the sense of self and subjective continuity.
        
        \bibliographystyle{plainnat}
        \bibliography{/home/donkarlo/Dropbox/repo/nd_shared_research/bibliography/bibliography.bib}
\end{document}